% =========================================================================
% GLOSARIO
% =========================================================================
\chapter*{Glosario}
\addcontentsline{toc}{chapter}{Glosario}
\begin{description}[leftmargin=3.2cm,style=nextline]
\item[Cartera vencida] Parte del crédito otorgado cuyos pagos presentan atrasos conforme a los criterios de la autoridad.
\item[Cartera vigente] Parte del crédito otorgado cuyos pagos están al corriente.
\item[Default] Evento binario definido en esta tesis como IMOR futuro superior a 10\%.
\item[dpi] Puntos por pulgada. Mide la resolución de una imagen; las figuras de esta tesis se generan a 350 puntos por pulgada.
\item[Ensamble] Combinación de múltiples modelos individuales para mejorar la capacidad predictiva.
\item[Estado de resultados] Estado financiero que muestra los ingresos, los gastos y el resultado de una institución durante un periodo.
\item[Estimaciones preventivas] Reservas contables que una institución constituye para cubrir pérdidas esperadas de su cartera de crédito.
\item[Fuga de información] Uso, durante el entrenamiento, de información que no estaría disponible al momento de hacer la predicción.
\item[Inflación] Aumento sostenido del nivel general de precios, medido en México mediante el INPC.
\item[Lookahead bias] Sesgo que ocurre cuando información no disponible en el momento $t$ se utiliza indebidamente para predecir eventos en $t$.
\item[Margen financiero] Diferencia entre los ingresos por intereses y los gastos por intereses de una institución.
\item[Panel no balanceado] Conjunto de datos panel donde el número de observaciones temporales difiere entre unidades transversales.
\item[Parquet] Formato de archivo columnar y comprimido para almacenar tablas de datos.
\item[Resultado neto] Utilidad o pérdida final de una institución después de gastos, estimaciones e impuestos.
\item[SOFIPO] Sociedad Financiera Popular, intermediario financiero regulado por la CNBV para segmentos populares.
\item[Tasa de interés real] Tasa de interés nominal descontada por la inflación.
\item[TimeSeriesSplit] Esquema de validación cruzada que respeta el orden cronológico de las observaciones.
\end{description}

